\subsection{ROLLE'S THEOREM}

PROPOSITION 1 (``Rolle's theorem''). Let f be a real function which is finite and continuous on a closed interval I = {[}a, b{]} (where a \textless{} b), has a derivative (finite or not) at every point of {]}a, b{[}, and is such that \(f(a) = f(b)\) . Then there exists a point c of {]}a, b{[} such that \(f'(c) = 0\) .

The proposition is evident if \(f\) is constant: if not, \(f\) takes, for example, values \(> f(a)\) , and so attains its least upper bound at a point \(c\) interior to I (Gen.~Top., IV, p.~359, th. 1). Since \(f\) has a relative maximum at this point we have \(f'(c) = 0\) (I, p.~20, prop. 7).

COROLLARY. Let \(f\) be a real function which is finite and continuous on \([a, b]\) (where \(a < b\) ), and has a derivative (finite or not) at every point. Then there exists a point \(c\) of \(]a\) , \(b[\) such that \(f(b) - f(a) = f'(c)(b - a)\) .

We need only apply prop. 1 to the function \(f(x) - \frac{f(b) - f(a)}{b - a} (x - a)\) .

This corollary signifies that there is a point \(\mathbf{M}_{c} = (c, f(c))\) on the graph C of f such that a \textless{} c \textless{} b and such that the tangent to C at this point is parallel to the line joining the points \(\mathbf{M}_{a} = (a, f(a))\) and \(\mathbf{M}_{b} = (b, f(b))\) .

